\section{Conclusion and Future Work}
\label{sec:conclusion}
Tool-level least privilege leaves a data agent's reach unchanged once its query tool is allowed. Task-scoped
entitlements close that gap for relational data: a broker infers the tables, columns and rows a task needs, the
database enforces that scope through ephemeral roles and views owned by the user's role, and a bounded widening
protocol absorbs inference errors. On \nTasks{} public tasks, TSE scopes reduced reachable columns by a factor of
\oerColReduction{} and reachable cells by a factor of \oerCellReduction{} under synthetic-policy entitlements
(\oerCellReductionSuff{} on tasks whose scope sufficed, and about \oerCellRedPostSonnet{} counting approved widening),
at an accuracy cost of \dExHaikuTse{} to \dExSonnetTse{} points, and lowered injection-driven exfiltration of
out-of-task data. Widening approval is the main residual channel: an injection that targets it succeeded in most
runs under automatic approval, so deployments should require human approval or deny widening for columns outside
the initial scope. The approach depends on scope-inference quality, which we measured directly, and on the
correctness of database enforcement, which we tested but which ultimately rests on the DBMS.

Future work includes multi-step tasks whose scope evolves with intermediate results; propagating row restrictions
along joins (semi-join reduction of related tables); scope inference for enterprise schemas with thousands of
columns; combining TSE with information-flow control on the agent side~\cite{costa2025fides,debenedetti2026camel},
so that in-scope data cannot flow to arbitrary sinks; and support for other database systems.
